\documentclass[10.5pt,letterpaper]{article}

\usepackage[margin=0.75in,top=0.6in,bottom=0.6in]{geometry}
\usepackage{titlesec}
\usepackage{enumitem}
\usepackage{hyperref}
\usepackage{fontspec}
\usepackage{xcolor}
\usepackage{tabularx}
\usepackage{microtype}

\setmainfont{Times New Roman}
\linespread{0.97}
\emergencystretch=2.5em

\definecolor{linkcolor}{RGB}{30,58,95}
\hypersetup{
  colorlinks=true,
  urlcolor=linkcolor,
  pdftitle={Fengyun Yu CV},
  pdfauthor={Fengyun Yu}
}

\pagestyle{empty}
\setlength{\parindent}{0pt}

\setlist{topsep=1pt,itemsep=0.5pt,parsep=0pt,partopsep=0pt}

% Section headings: bold, full-width rule below
\titleformat{\section}{\large\bfseries}{}{0pt}{}[\vspace{-6pt}\hrulefill]
\titlespacing{\section}{0pt}{8pt}{4pt}

\newcommand{\entry}[4]{%
  \textbf{#1} \hfill #2\\
  #3 \hfill \textit{#4}\\[1pt]
}

\newcommand{\subentrytitle}[2]{%
  \textbf{#1} \hfill #2\\[1pt]
}

\begin{document}

\begin{center}
  {\LARGE \bfseries \MakeUppercase{Fengyun Yu}}\\[4pt]
  Interdisciplinary Center for Scientific Computing\\
  Universit\"at Heidelberg, Germany\\
  e: \href{mailto:yufengyun17@gmail.com}{yufengyun17@gmail.com}\\
  \href{https://scholar.google.com/citations?user=7bPHV-AAAAAJ&hl=en}{Google Scholar} \textbar\
  \href{https://www.researchgate.net/profile/Yu-Feng-Yun?ev=hdr_xprf}{Research Gate} \textbar\
  \href{https://github.com/Finley-Maple}{GitHub} \textbar\
  \href{https://www.linkedin.com/in/fengyun-yu}{LinkedIn}
\end{center}

\section*{Education}

\entry{Heidelberg University}{Heidelberg, Germany}{Master's degree in Scientific Computing}{2023.4 -- 2026.4}
\begin{itemize}[leftmargin=18pt,itemsep=1pt,topsep=2pt]
  \item Graduated with overall grade \textbf{1.3/1.0} (German scale, ``sehr gut''; 1.0 = highest)
  \item Master's thesis: \textit{Digesting Cell Segmentation in Spatial Transcriptomics via Graph Representation Learning} (grade 1.3), advised by Elyas Heidari and Prof. Moritz Gerstung
  \item Application area: Economics (Grade: 1.35)
  \item Relevant coursework: Generative neural network (1.0), Theory of deep learning (1.0), Infinite dimension optimization (1.3), NLP with Transformer (1.3), Mathematical machine learning (1.3)
\end{itemize}

\entry{Tsinghua University}{Beijing, China}{Bachelor's degree in Industrial Engineering}{2017.7 -- 2021.6}
\begin{itemize}[leftmargin=18pt,itemsep=1pt,topsep=2pt]
  \item Average grade: 3.67/4.00 (equivalent to 1.3 in German scale as per APS); minor in Statistics
  \item Honor: Evergrande Scholarship (top 5\% of the cohort)
  \item Exchange semester: KTH Royal Institute of Technology, Stockholm, Sweden (2019.8 -- 2020.1)
\end{itemize}

\section*{Research Experience and Interest}

\textbf{Research Interest:}
\begin{enumerate}[leftmargin=18pt,itemsep=1pt,topsep=2pt]
  \item Trustworthy, interpretable machine learning for health and biomedicine, where accuracy alone is not evidence of meaning; currently geometric deep learning for spatial omics data
  \item Combination of causal inference and machine learning in medical data, with a focus on identifying causal relationships in observational studies and clinical outcomes
\end{enumerate}
\textbf{Additional interests:} Empirical health economics studies, numerical optimization.

\section*{Bibliography}

\textbf{Peer-reviewed articles}
\begin{enumerate}[leftmargin=18pt,itemsep=3pt,topsep=2pt]
  \item Wang C, \textbf{Yu F} (Co-first), Cao Z, et al. Exploring COPD Patient Clusters and Associations with Health-Related Quality of Life Using A Machine Learning Approach: A Nationwide Cross-Sectional Study[J]. \textit{Engineering}, 2025. \url{https://doi.org/10.1016/j.eng.2025.05.005}
  \item \textbf{Yu F}, Jiao L, Chen Q, Wang Q, De Allegri M, et al. (2024) Preferences regarding COVID-19 vaccination among 12,000 adults in China: A cross-sectional discrete choice experiment. \textit{PLOS Global Public Health} 4(7): e0003387. \url{https://doi.org/10.1371/journal.pgph.0003387}
  \item \textbf{Yu F}, Geldsetzer P, Meierkord A, Yang J, Chen Q, Jiao L, Abou-Arraj NE, Pan A, Wang C, B\"arnighausen T, Chen S. Knowledge About COVID-19 Among Adults in China: Cross-sectional Online Survey. \textit{J Med Internet Res}. 2021 Apr 29;23(4):e26940. \url{https://doi.org/10.2196/26940}. PMID: 33844637; PMCID: PMC8086781.
  \item Chen S, Kuhn M, Prettner K, \textbf{Yu F}, Yang T, B\"arnighausen T, Bloom DE, Wang C. The global economic burden of chronic obstructive pulmonary disease for 204 countries and territories in 2020-50: a health-augmented macroeconomic modelling study. \textit{Lancet Glob Health}. 2023 Aug;11(8):e1183-e1193. \url{https://doi.org/10.1016/S2214-109X(23)00217-6}
\end{enumerate}
Additional papers: refer to \href{https://scholar.google.com/citations?user=7bPHV-AAAAAJ&hl=en}{Google Scholar}

\section*{Major Research Experience}

\subentrytitle{German Cancer Research Center (DKFZ), Lab of AI in Oncology}{Heidelberg, Germany}
\subentrytitle{Interpretability and benchmarking of Segger (E. Heidari), a graph neural network for cell segmentation --- Master's thesis}{2025.3 -- 2026.3}
\begin{itemize}[leftmargin=18pt,itemsep=1pt,topsep=2pt]
  \item Built an interpretability framework for Segger: extracted and analyzed its attention weights layer by layer and visualized transcript--cell and gene--gene embedding manifolds (UMAP), showing the attention space captures real gene--gene and gene--cell interaction patterns; a white-box route to unsupervised discovery of spatially variable genes
  \item Studied Aligned Segger, which adds an alignment loss on mutually exclusive gene sets to sharpen cell boundaries and reduce ``leaky RNA''; showed the alignment loss ($\lambda=0.5$), more than model width, reshapes the embedding toward known biology, and that high raw scores (AUROC 0.979, F1 0.910) can decouple from biological meaningfulness
  \item Benchmarked Segger against Proseg, Bering, Baysor, and BIDCell on 8 cancer datasets; only Segger and Proseg scale beyond 100 million transcripts, while the others fail with out-of-memory or timeouts
\end{itemize}

\subentrytitle{Heidelberg Institute of Global Health (HIGH)}{Heidelberg, Germany}
\subentrytitle{Time-encoding benchmark for mortality risk prediction}{2026.1 -- 2026.6}
\begin{itemize}[leftmargin=18pt,itemsep=1pt,topsep=2pt]
  \item Systematically compared four time-encoding strategies (decoupled sinusoidal, learnable positional embeddings, text serialization, LLM encoding) for mortality risk prediction on UK Biobank (n=124,158)
  \item Best-performing approach improved C-index by +0.008 over baseline (bootstrap p\textless0.001)
\end{itemize}

\subentrytitle{Universit\"at Heidelberg}{Heidelberg, Germany}
\subentrytitle{Deep Generative Models: generative assignment flows for representing discrete data}{2024.9 -- 2025.3}
\begin{itemize}[leftmargin=18pt,itemsep=1pt,topsep=2pt]
  \item Introduced a novel generative model for the representation of joint probability distributions of discrete random variables by projecting high-dimensional flow matching
\end{itemize}

\subentrytitle{Chinese Academy of Medical Sciences \& Peking Union Medical College (CAMS\&PUMC)}{Beijing, China}
\subentrytitle{Exploring COPD patient clusters and associations with health-related quality of life}{2023.12 -- 2025.3}
\begin{itemize}[leftmargin=18pt,itemsep=1pt,topsep=2pt]
  \item Co-first author; published in \textbf{Engineering} (2025)
  \item Led the clustering analysis (k-means, hierarchical clustering, latent class analysis) on a nationwide COPD cohort (n\textgreater10,000); identified four clinical phenotypes and quantified their association with health-related quality of life (EQ-5D)
\end{itemize}

\subentrytitle{Chinese Academy of Medical Sciences \& Peking Union Medical College (CAMS\&PUMC)}{Beijing, China}
\subentrytitle{Preferences and willingness to pay for COVID-19 vaccine among adults in China}{2020.12 -- 2024.6}
\begin{itemize}[leftmargin=18pt,itemsep=1pt,topsep=2pt]
  \item First author; published in \textbf{PLOS Global Public Health} (2024)
  \item Quantified vaccine preference using a mixed logit model on a nationally representative discrete choice experiment sample (n=12,000)
  \item Contributed to 5+ peer-reviewed publications (incl. Lancet Global Health, JMIR) during a research assistantship at CAMS\&PUMC, 2021--2023
\end{itemize}

\section*{Skills and Interests}

\begin{itemize}[leftmargin=18pt,itemsep=1pt,topsep=2pt]
  \item \textbf{Technical:} Python (pandas, numpy, scikit-learn, PyTorch, Streamlit, Plotly), R, SQL, MATLAB, C/C++, TypeScript/JavaScript, LaTeX
  \item \textbf{Methods:} Graph neural networks, mixed-effects models, discrete choice experiments / mixed logit, survival analysis, Bayesian statistics, clustering, optimal transport, reinforcement learning
  \item \textbf{Languages:} Fluent in English, Chinese (native), intermediate German (B1)
  \item \textbf{Interests:} Marathon running, badminton, swimming
\end{itemize}

\end{document}
